% !TeX root = ../papers/nc-p3.tex
\section{片側分岐を結ぶ有限な有理曲線列と次数公式}
\label{sec:finite-branch-corridors}
\begin{theorem}\label{thm:finite-branch-corridors}
全射影像が高々一次元の標準正則窓を考える。後続直線を持つ点 $p$ を一つ固定する。
その直線のどの点でも、その辺における前点は $p$ 一つである。
また $p$ から出発する有理曲線列は、有理的な一般点の後続を一意に取ることにより
\[
 x_0=p,\quad x_1,\ldots,x_m:\mathbb P^1\longrightarrow\mathbb P^3,
 \quad x_{m+1}=z
\]
となる。中間の射は全て像への双有理射で、$x_1,x_m$ の像は直線、$m\ge2$ である。
$z$ は前点の直線を持ち、反対向きの曲線列はこの列を逆にたどる。
従って四周期で同定した順側分岐中心と逆側分岐中心はこの列で一対一に対応する。
\end{theorem}
\begin{proof}
前節により $p$ の逆支持は三次元である。後続直線上の二項変数を $x$ とすると、
$R_1$ のその直線への射影は三次元である。その基底を $K_j(x)\in V_2$ として
$w(x)=\sum_{j=1}^3 K_j(x)\otimes A_j$ と書く（位置を巡回して読む）。
一般点での $K_j(x)$ は独立で、$A_j$ の位置0支持の和は $\ker p$ である。
これはその点の逆支持が、可動中心の不在により一般には三次元だからである。
従って $K_j(x)$ が独立な全ての点で逆支持は三次元である。
従属な点の順支持は二次元であり、定理\ref{thm:no-two-sided-centers}によってやはり逆支持は三次元である。

一般点に先行する点の逆支持が3のとき、直交不等式から一般点の順支持は3以下である。
動く曲線の全点が階数2にはならないので一般に3で、一意な後続への有理写像を得る。
これは射影直線上で射に延びる。像が動けば、その一般点も一意な前点を持つので
標数零より曲線間の写像は双有理である。像が一定なら直前の曲線はその前点の直線全体となる。
その正規化パラメータはその直線との同型である。
いずれの動く像も全射影像の既約成分であり、位置を法4で見れば有限個しかない。
成分が二度現れると、一意な一般前点を順に復元して最初の直線を二度目にも得るが、
初回には前点が一定、二度目には変動するので矛盾する。従って有限回で一定像に達する。
逆の構成も一意なので最後の対応は全単射である。$m=1$ の排除は以下の次数式で示す。
\end{proof}

\begin{proposition}\label{prop:branch-corridor-degree}
$d_i=\deg x_i^*\mathcal O(1)$ とする。$x_{-1},x_{m+2}$ を両端の外の一意な定点とし、
動かない位置の次数を0とする。$1\le i\le m$ で逆・順の階数3評価の三次小行列式の
共通零因子の次数をそれぞれ $b_i,f_i$ とする。$b_0=f_{m+1}=0$ と置く。このとき
\[
 d_{i+2}-3d_{i+1}+3d_i-d_{i-1}= b_i-f_{i+1}\qquad(0\le i\le m).
\]
$b_1=f_m=0$ であり、
\[
 \sum b_i=\sum f_i\ge2,\qquad
 \sum_{i=1}^m\big(i b_i-(i-1)f_i\big)=m+2 .
\]
特に順側分岐中心が一つでもあれば、順側分岐中心を通る実際の周期点列がある。
\end{proposition}
\begin{proof}
辺 $x_i,x_{i+1}$ に沿い、$w$ による隣接六空間の完全対合を使う。
$x_i$ を評価して得る逆側の三空間は、列として $V_{i-1}^*$ から定数六空間へ入る
$d_i$ 次の行列の像であり、有理的な核直線は $x_{i-1}$ である。
従って飽和した像束の双対行列式の次数は $3d_i-d_{i-1}-b_i$ である。
実際 $V_{i-1}^*\otimes\mathcal O(-d_i)$ からの写像の核は
$\mathcal O(-d_i-d_{i-1})$、飽和で増える次数は局所単因子の和、すなわち三次小行列式の共通零因子の長さである。
隣の $x_{i+1}$ の順側評価でも次数は $3d_{i+1}-d_{i+2}-f_{i+1}$ である。
一般点で二つは直交補なので、その飽和三空間へのGrassmann写像は双対同型で対応し、Plücker次数が等しい。
最初・最後の定点で用いた側は階数3一定なので同じ計算が適用できる。
最初の直線上の逆支持は全点で3であるから $b_1=0$、逆に $f_m=0$ である。
$m=1$ なら最初の次数式から $f_1=3$ となり矛盾する。

次数式を加えれば、二端で二階差分がともに1なので総和は零である。
$i$ 倍して加えると、端の傾きが $+1,-1$ であることから表示した $m+2$ となる。
共通総和が1以下なら左辺は高々 $m$ であり矛盾する。
従って各列の内部で少なくとも一つ順側分岐中心に達する。実際のパラメータ値でその点まで進んでから
その中心の後続直線で同じことを繰り返せる。四位置合わせて中心は有限個なので、
ある中心を再び訪れ、間の有限経路を繰り返せば周期点列を得る。
\end{proof}
